\section{Parámetros del vehículo de referencia y del catálogo}
\label{sec:parametros}

\subsection{Por qué el PAE-X450}

Un modelo con ecuaciones correctas y parámetros inventados no sirve para
sintonizar nada. La 2.ª edición armaba su vehículo con números de distintas
tesis, y la \cref{sec:proposito} mostró que ni siquiera dentro de una misma
tesis eran coherentes. La 3.ª edición adopta como vehículo por defecto un
cuadricóptero en X de clase 450~mm y 1.4~kg (motores sin escobillas, ESC y
hélices de unos 10~pulgadas: el ``dron mediano normal'' de laboratorio) cuyos
parámetros vienen \emph{todos} de una misma fuente: el modelo por defecto de
CopterSim, el simulador del grupo de Quan Quan en la Universidad de Beihang,
base de la plataforma RflySim y de los libros \cite{qq17,qq20}.

La fuente primaria es el archivo de inicialización del modelo Simulink,
\cod{Init.m}, del repositorio \cod{github.com/XunhuaDai/CopterSim}, commit
\cod{0d14487f} (25 de julio de 2019) \cite{cs19}, en el que cada parámetro
viene comentado. Los valores de la \cref{tab:pae} se cotejaron contra ese
archivo el 4 de octubre de 2026 y coinciden dígito a dígito.

\begin{table}[H]
\centering
\small
\caption{Parámetros del PAE-X450 (\cod{vehicles.PAE\_X450}). Fuente: \cod{Init.m}
de CopterSim \cite{cs19}, modelo por defecto de \cite{qq20,qq17}, salvo donde
se indica.}
\label{tab:pae}
\setlength{\tabcolsep}{4pt}
\begin{tabularx}{\textwidth}{>{\raggedright\arraybackslash}p{3.9cm}l>{\raggedright\arraybackslash}p{2.6cm}l>{\raggedright\arraybackslash}X}
\toprule
Parámetro & Símbolo & Valor & Unidad & Nombre en \cod{Init.m} / nota \\
\midrule
Masa & $m$ & 1.4 & kg & \cod{uavMass} \\
Brazo (centro a eje del motor) & $\ell$ & 0.225 & m & \cod{uavR} \\
Inercias & $J_{xx},J_{yy},J_{zz}$ & 0.0241, 0.0239, 0.0386 & kg·m$^2$ & \cod{uavJxx/yy/zz}; productos de inercia nulos \\
Coeficiente de empuje & $C_T$ & $1.105\times10^{-5}$ & N/(rad/s)$^2$ & \cod{rotorCt} \\
Coeficiente de par & $C_M$ & $1.779\times10^{-7}$ & N·m/(rad/s)$^2$ & \cod{rotorCm} \\
Inercia rotor + hélice & $J_m$ & $1.287\times10^{-4}$ & kg·m$^2$ & \cod{motorJm} \\
Pendiente del ESC & $C_R$ & 1148 & rad/s & \cod{motorCr} \\
Ordenada del ESC & $\omega_b$ & $-141.4$ & rad/s & \cod{motorWb} \\
Constante de tiempo del motor & $T_m$ & 0.02 & s & \cod{motorT} \\
Zona muerta del acelerador & $\sigma_{dz}$ & 0.05 & --- & \cod{motorMinThr} \\
Arrastre horizontal & $C_{d,xy}$ & 0.055 & N/(m/s)$^2$ & \cod{uavCd} \\
Arrastre vertical & $C_{d,z}$ & 0.055 & N/(m/s)$^2$ & \textbf{supuesto}: igual a $C_{d,xy}$ \\
Amortiguamiento rotacional & $C_{Cm}$ & 0.0035, 0.0039, 0.0034 & N·m/(rad/s)$^2$ & \cod{uavCCm} \\
Acelerador de hover declarado & --- & 0.609 & --- & \cod{BattHoverThr}; solo para comprobar \\
\bottomrule
\end{tabularx}
\end{table}

\begin{supuesto}[Lo que no viene de CopterSim]
El arrastre vertical; el voltaje constante (CopterSim trae un modelo de
batería que aquí no se implementa); y el punto de aplicación del arrastre:
\cod{Init.m} declara una distancia vertical de 0.12~m entre el centro
aerodinámico y el centro de masa (\cod{uavDearo}), que produciría un momento
de cabeceo proporcional al arrastre. El simulador aplica el arrastre en el
centro de masa y no modela ese momento (\cref{sec:limitaciones}).
\end{supuesto}

\subsection{Chequeos de consistencia}

Un conjunto de parámetros tomado de una fuente puede tener erratas. La
forma de detectarlas es cruzar cantidades que la fuente da por separado y
que la física obliga a que coincidan. El bloque A de
\cod{model\_verification.py} hace cinco cruces.

\paragraph{Velocidad y acelerador de hover.} El rotor sostiene el hover
cuando $4C_T\Omega_h^2=mg$:
\begin{equation}
\Omega_h=\sqrt{\frac{mg}{4C_T}}=\val{wh}~\text{rad/s}\ (\val{whrpm}~\text{rpm}),\qquad
\sigma_h=\frac{\Omega_h-\omega_b}{C_R}=\val{sigmah}.
\label{eq:hover}
\end{equation}
\cod{Init.m} declara por separado un acelerador de hover de
\val{sigmahd}. Que el valor calculado desde masa, $C_T$ y la recta del ESC
coincida con el declarado a la tercera cifra indica que propulsor y masa son
coherentes entre sí: es exactamente la prueba que el F450 de Jiřinec no pasa.

\paragraph{Coeficiente de empuje adimensional.} La literatura de hélices
reporta el coeficiente estático $C_T'=T/(\rho n^2D^4)$, con $n$ en rev/s. Con
$\Omega=2\pi n$, $C_T'=4\pi^2C_T/(\rho D^4)$. Suponiendo la hélice de
9.4~pulgadas (0.239~m) que monta un F450 (supuesto: \cod{Init.m} no declara
el diámetro) y $\rho=1.225$~kg/m$^3$, resulta $C_T'=\val{CTad}$ y
$C_P'=2\pi C_Q'=\val{CPad}$. Son valores del orden de los que publican las
bases de ensayos estáticos para hélices de dos palas de ese tamaño (UIUC
Propeller Data Site \cite{uiuc}), lo que descarta un error de unidades en
$C_T$ y $C_M$, por ejemplo una confusión entre rpm y rad/s, que cambiaría el
coeficiente en un factor $(60/2\pi)^2\approx91$.

\paragraph{Relación par/empuje.} $C_M/C_T=\val{CmCt}$~m, típica de
hélices de 9--10~pulgadas; de ella sale la autoridad de guiñada.

\paragraph{Empuje máximo.} Con $\sigma=1$, $\Omega_{max}=C_R+\omega_b=\val{wmax}$~rad/s
(\val{wmaxrpm}~rpm) y cada rotor da $C_T\Omega_{max}^2=\val{Tmax}$~N
(\val{Tmaxkgf}~kgf). La relación empuje/peso es $\val{TW}$, normal para un
F450 con batería 4S y suficiente para maniobrar con margen.

\paragraph{Frecuencia de paso de pala.} Con dos palas, la vibración por
desbalance aparece a $f_{bp}=2\Omega/(2\pi)$: \val{fbp}~Hz en hover y
\val{fbpmax}~Hz a acelerador máximo. Este número gobierna el diseño de los
filtros (\cref{sec:sensores,sec:filtros}).

\subsection{Otros vehículos del catálogo}

El catálogo (\cod{VEHICLES}) conserva cuatro vehículos más, útiles para
mostrar que el método de sintonía se adapta solo cuando cambia la física. La
\cref{tab:catalogo} resume sus parámetros y las cantidades derivadas
calculadas por el mismo código.

\begin{table}[H]
\centering
\small
\setlength{\tabcolsep}{5pt}
\caption{Catálogo de vehículos (valores y derivados calculados por \cod{vehicles.py}).}
\label{tab:catalogo}
\begin{tabular}{lccccc}
\toprule
 & \val{catHead} \\
\midrule
\val{catT}
\bottomrule
\end{tabular}
\end{table}

Las fuentes y advertencias de cada uno están en sus campos \cod{source} y
\cod{caveats}:
\begin{itemize}
\item \textbf{Jiřinec (+)} \cite{jirinec11}: inercias (ec.~4.7), $b$ (ec.~4.9),
$d$ (p.~28), motor $0.7/(0.1s+1)$ (ec.~4.6), brazo 0.214~m. Corregido a
1.321~kg y bastidor + (\cref{sec:proposito}). ESC lineal sin
desplazamiento y empuje/peso 3.0 son supuestos; el arrastre lineal
0.25~kg/s viene de Luukkonen.
\item \textbf{Luukkonen} \cite{luukkonen11} (tabla~1, que toma los valores de
Tayebi y McGilvray): rotores instantáneos en la fuente, por lo que
$T_m=0.05$~s es supuesto.
\item \textbf{NUS KayLion} \cite{hb11}: el único con todo medido en banco
(ecs.~11.18--11.27), incluida $T_m=0.0821$~s. Empuje/peso 2.0 deducido de
su banco.
\item \textbf{AscTec Hummingbird} \cite{hb16} (tablas~16.1--16.2): sus
coeficientes venían en N/rpm$^2$; convertidos dividiendo por
$(2\pi/60)^2$ dan $k_F=\val{hbkF}$ y $k_M=\val{hbkM}$ en unidades SI.
\end{itemize}
En los cuatro, el ESC se modela como recta sin desplazamiento
($\omega_b=0$) con $C_R$ elegido para la relación empuje/peso indicada, lo
que fija $\sigma_h=1/\sqrt{T/W}$. Para estos vehículos el acelerador de
hover es una consecuencia de ese supuesto, no un dato.
